\section{Canonical trace compatibility between distinct actual power charts}
\label{sec:kn-nadel-kummer-cross-chart}

Fix the original coefficient algebra $R\to S$ and an integer $n$.
For $i=1,2$, let $a_i:R[x_1,\ldots,x_n]\to S$ be an actual
$R$-algebra homomorphism, $d_i:\{1,\ldots,n\}\to\mathbf N$, and
let $C_i=S\otimes_{(R[x],\,x_j\mapsto t_j^{d_i(j)+1})}R[t]$ be
the literal tensor algebra.  Write $\ell_i:S\to C_i$ for its actual
left inclusion.  The two degrees, chart maps, tensor algebras, and
scalar actions are retained separately throughout.

\begin{lemma}[The actual common original base top-form module]
\label{lem:kn-nadel-cross-chart-original-coefficients}
\leanok
The original algebra condition on each $a_i$ internally identifies
the chart's genuine base top-form module with the same original
$\omega_S=\bigwedge^n_S\Omega_{S/R}$.  This identification fixes
every actual ordered differential wedge, not only the chosen chart
generator.  In particular it supplies the actual comparison between
the two base top-form modules without a canonical-comparison input.
\end{lemma}
\begin{proof}
The whole small coefficient triangle is
$a_i\circ C=\operatorname{algebraMap}_{R,S}$ by the actual native
algebra homomorphism law.  The induced algebra structure on $S$
therefore equals the original one.  Transport the actual native
K\"ahler module and its top exterior module along this equality.
After that actual coefficient identification the map is the identity,
so it fixes every wedge of actual differentials.  Composing the first
identification with the inverse of the second gives the claimed
whole base comparison.
\uses{thm:kn-nadel-etale-top-frame-base-change}
\end{proof}

\begin{theorem}[The genuine differential square for an actual cross-chart algebra isomorphism]
\label{thm:kn-nadel-cross-chart-exterior-square}
\leanok
Suppose $a_1$ is formally etale.  Let $e:C_1\simeq_S C_2$ be an actual coefficient-preserving native
algebra isomorphism.  It induces an actual original $R$-algebra
isomorphism and hence a genuine K\"ahler exterior map
$E:\omega_{C_1}\to\omega_{C_2}$.  After identifying the two actual
base top-form modules with the original $\omega_S$, the genuine
differential pullbacks satisfy
\[
 E(P_1(w))=P_2(w)\qquad(w\in\omega_S).
\]
The square is proved from the actual coefficient homomorphisms;
no differential or canonical comparison diagram is an input.
\end{theorem}
\begin{proof}
The whole tensor coefficient square identifies the original
$R\to C_i$ map with $R\to S\xrightarrow{\ell_i} C_i$.
Because the actual $S$-algebra isomorphism fixes these coefficient
images, it preserves the actual original $R$-structures.  Apply
the native algebra K\"ahler morphism and its exterior universal
property to construct $E$.

On any differential wedge of actual functions $s_j\in S$, the
left composite is the wedge of $D_{C_2}(e(\ell_1(s_j)))$; the
right composite is the wedge of $D_{C_2}(\ell_2(s_j))$.  The actual
coefficient-preserving law gives equality term by term.  The
previous lemma identifies the base differential wedges exactly.
The first genuine etale chart's top-form basis, together with
actual $S$-linearity and semilinearity through the same coefficient
maps, extends this equality to every actual base top form.
\uses{lem:kn-nadel-cross-chart-original-coefficients,thm:kn-nadel-etale-all-topforms-trace-normalization}
\end{proof}

\begin{theorem}[The whole canonical trace is independent of the distinct power-chart presentation]
\label{thm:kn-nadel-cross-chart-canonical-counit}
\leanok
Suppose $R$ is a characteristic-zero domain, both $a_i$ are formally
etale, and $a_1$ is flat.  For the actual cross-chart algebra
isomorphism above, let $\tau_i:\omega_{C_i}\to\omega_S$ be the
already constructed actual canonical counits, followed by the
internally produced common-base identifications.  Then
\[
 \tau_2\circ E=\tau_1
\]
as whole actual $S$-linear maps.  The two degree functions and
presentations need not agree.  No trace, canonical, differential,
or vanishing comparison is supplied as a premise.
\end{theorem}
\begin{proof}
The genuine exterior map $E$ is $S$-linear after restriction through
the two actual inclusions, since $e(\ell_1(s))=\ell_2(s)$.  Its
composite with the second actual counit and the proved base
identifications is thus an actual candidate counit for the first
chart.  For every actual $c\in C_1$ and $w\in\omega_S$, the
preceding differential square and the proved all-top-form
normalization on the second chart give
\[
 (\tau_2 E)(cP_1(w))
   =\tau_2(e(c)P_2(w))
   =\operatorname{Tr}_{C_2/S}(e(c))w
   =\operatorname{Tr}_{C_1/S}(c)w.
\]
The last equality is native trace invariance under the actual
$S$-algebra isomorphism, with each original scalar action retained.
The whole-counit uniqueness theorem for the first actual flat chart
therefore identifies this composite with its internally normalized
canonical counit.  Transport through the actual common-base
identification gives the displayed whole-map equality.
\uses{thm:kn-nadel-cross-chart-exterior-square,thm:kn-nadel-etale-all-topforms-trace-normalization,thm:kn-nadel-etale-canonical-counit-uniqueness}
\end{proof}

The actual coefficient-preserving algebra isomorphism is geometric
ring data; no desired differential compatibility is part of it.
Producing these ring identifications for the original whole-cover
overlaps and subsequently gluing the canonical sheaf maps remain
separate geometric tasks.  This theorem does not assert an ordinary
canonical line on a singular normalized cyclic cover or completed
Nadel vanishing.
